
Como ya habíamos mencionado, antes de quitar el acordonamiento, debe vigilarse que la fuente se encuentre dentro del contenedor. Si no se encuentra en la posición deseada, se deberá:

\begin{itemize}
  \item[a)] Mantener vigilada la zona restringida, a fin de que nadie penetre en la misma. Avisar a las personas de la situación y alejarlas.
  \item[b)] Uno de los técnicos se mantendrá en el área de menor exposición (con ayuda del detector) vigilando el equipo, mientras que el otro solicita ayuda al encargado de la seguridad radiológica, para que éste acuda con los accesorios para recuperar las fuentes (si no cuentan con ellos) y colabore en la recuperación.
  \item[c)] Localizar la posición de la fuente con el equipo detector con binoculares.
\end{itemize}

Los accesorios mínimos con que se deberá contar para la recuperación de fuentes son: contenedor de recuperación, caña con imán en la punta y embudo. (Fig. 1).

El contenedor de recuperación deberá tener un espesor suficiente para atenuar la rapidez de exposición a niveles aceptables, tomando como base la fuente de más alta actividad que pudiera recuperarse. Sus dimensiones internas deberán permitir que la fuente se introduzca hasta las 3/4 partes de su longitud.

La caña deberá tener una longitud mínima de 4 metros, lo que permitirá recuperar una fuente de 100 Ci de iridio-192 recibiendo una exposición de 50 mR si la operación dura 1 minuto.
